\subsubsection{Função de transição}

Como abordado na seção \ref{fsm}, a dinâmica do sistema é descrita por uma função de transição $\delta$, a qual mapeia o estado atual do sistema para seu estado sucessor a cada evento de clock.

Seja $Q_t = (H_t, M_t, S_t)$ o estado do sistema no instante $t$. A evolução temporal é definida por:

\[
Q_{t+1} = \delta(Q_t)
\]

No contexto deste projeto, a função de transição é induzida por uma operação de incremento ou decremento, definida pelo parâmetro $\sigma_t \in \Sigma$, que indica a operação de incremento ou decremento.

\[
\delta(H_t, M_t, S_t; \sigma_t) = (H_{t+1}, M_{t+1}, S_{t+1})
\]

Além disso, para esta estruturação dos estados, $\delta$ foi escolhida de forma na qual a atualização do estado não ocorre de forma uniforme sobre todos os componentes, mas segue uma estrutura hierárquica com propagação condicional entre dígitos.
Dessa forma, introduz-se um conjunto de sinais de enable $e$.

\begin{align*}
e_S &= 1 \\
e_M &= c_S \\
e_H &= c_M \land c_S
\end{align*}

onde $c_S$ e $c_M$ representam os sinais de carry/borrow gerados pelos contadores de segundos e minutos, respectivamente.

\begin{align*}
    (S_{t+1}, c_S) &= \delta_S(S_t, \sigma_t, e_S) \\
    (M_{t+1}, c_M) &= \delta_M(M_t, \sigma_t, e_M) \\
    (H_{t+1}) &= \delta_H(H_t, \sigma_t, e_H)
\end{align*}

Alternativamente: 

\begin{align*}
    \delta_S(S_t; \sigma) \rightarrow (S_{t+1}, c_S)\\
    \delta_M(M_t; \sigma) \xrightarrow{c_S} (M_{t+1}, c_M) \\
    \delta_H(H_t; \sigma) \xrightarrow{c_S \land s_M} (H_{t+1})
\end{align*}

A função global $\delta$ é, portanto, implementada por meio da composição das funções $\delta_S$, $\delta_M$ e $\delta_H$, com propagação de sinais de carry/borrow entre os módulos.

\[
\delta(Q_t) = \delta_h \circ \delta_m \circ \delta_s (Q_t)
\]

De forma análoga, cada contador é decomposto em unidades e dezenas, seguindo a mesma lógica hierárquica de habilitação interna:

\begin{align*}
(X_{u,t+1}, c_{int}) &= \delta_X^{(u)}(X_{u,t}, e_X) \\
(X_{d,t+1}, c_X) &= \delta_X^{(d)}(X_{d,t}, c_{int}, e_X)
\end{align*}

onde:
\begin{itemize}
    \item $X_u$ representa a unidade do dígito;
    \item $X_d$ representa a dezena;
    \item $c_{int}$ é o carry/borrow interno;
    \item $e_X$ é o sinal de habilitação do módulo.
    \item $c_X$ é dado por $(c_{int} \land e_X)$
\end{itemize}

Dessa forma, cada transição de estado pode ser implementada como um contador modular independente, cuja lógica de transição depende apenas de um subconjunto restrito de variáveis e de sinais de propagação provenientes do nível imediatamente inferior.